\section{Threats to Validity}
\textbf{Label noise.} KW-0929 annotation is literature-derived and not
uniformly experimentally verified; negative labels are absence of annotation,
not measured inactivity. Both noise sources bias metrics downward, so our
numbers are, if anything, conservative.

\textbf{Corpus shift.} UniProt and APD6 overlap through shared primary
literature. The 10-mer decontamination removed 1{,}795 of 3{,}207 usable APD6
entries, indicating substantial overlap; the retained 1{,}412 are the hard,
non-memorized cases, and the external AUROC of $0.926$ is measured on those.

\textbf{Split strictness.} Shared-10-mer clustering is stricter than typical
identity thresholds for short peptides and may split some true homolog pairs
that lack an exact decamer. A sensitive-alignment-based clustering (e.g.\
MMseqs2) is a planned upgrade; we expect similar conclusions at slightly
higher absolute numbers.

\textbf{Multiple comparisons.} Four models on three evaluations share one
test split; we report all planned comparisons with intervals rather than
selecting the best, and the permutation calibration ($p=0.0020$ at $B=500$)
is a resolution limit, not a sharp $p$-value.

\textbf{Sandbox constraints.} Truncating training at 30{,}000 sequences and 8
epochs caps absolute performance; the internal/external agreement suggests
the cap costs little, but the claim is about the pipeline, not a ceiling.
